% from main.tex
%%%%%%%%%%  scan idx 409  |  folio 396  %%%%%%%%%%
(\(c_i\) et \(c_i^{(j)}\)
étant effectués du poids \(i\)), et ces \(n\) relations engendrent l'idéal des relations entre les \(c_i^{(j)}\).
\end{enumerate}

\textbf{Corollaire 4.4.} Les éléments de \(C\) de la forme~:
\[
  \xi_1^{\alpha_1}\cdots\xi_{n-1}^{\alpha_{n-1}},
  \qquad 0\le\alpha_i\le n-i,
\]
forment une base de \(C\) considéré comme module sur \(A\).

Pour la démonstration du lemme et du corollaire, nous renvoyons à [2], exposé 4, pages 4.19, auquel ils sont empruntés.

\subsection*{4.5}
Soient \(S\) un schéma quasi-compact, \(E\) un \(\mathcal O_S\)-Module localement libre de rang \(n\), \(\pi=(p_1,\ldots,p_k)\) une suite d'entiers \(\geq1\), de somme \(n\). Posons \(X=D_\pi(E)\) ; \(X\) est muni d'une suite de quotients canonique \(L_i\) ; on pose \(L_0=0\), \(L_k=f^*(E)\), \(f\) étant le morphisme structural \(X\to S\). Pour \(1\le i\le k\) on note \(F_i\) le noyau de \(L_i\to L_{i-1}\). Soit \(K^{\bullet}(X)\) le groupe de Grothendieck des faisceaux localement libres de type fini sur \(X\) (resp. des complexes parfaits). On note \(F_j\) la classe de \(F_j\), et \(\lambda_i^{(j)}\) la classe de \(\Lambda^i(F_j)\). Dans le cas où \(K^{\bullet}(X)\) est défini en termes de complexes parfaits, on donne encore un sens évident aux séries \(\lambda_t(E)\), \(\lambda_t(F_j)\), où \(E\) est la classe de \(E\). Enfin, on remarque que~:
\[
  \sum_{j=1}^{k}F_j=E.
\]

\textbf{Théorème 4.6.} Avec les hypothèses et les notations de 4.5, la \(K^{\bullet}(S)\)-algèbre \(K^{\bullet}(X)\) est engendrée par les \(\lambda_i^{(j)}\); l'idéal des relations entre les \(\lambda_i^{(j)}\) est engendré par les relations déduites de~:
\[
  \lambda_t(E)=\prod_{j=1}^{k}\lambda_t(F_j).
\]
